% ==========================================================================
% ======================== Dokumentace k PC 1.0 ;-) ========================
% ==========================================================================

% Ověřeno, prošlo bez jediné připomínky - při dodržení následujících pokynů:
% - SPRÁVNĚ vyplnit titulní stranu (!!!)
% - Nadpisy používat v pořadí (zanoření):
% 		- chapter (celá kapitola)
%		- section
% 		- subsection, subsubsection ... 
% - Pro pomlčky V TEXTU používat -- (doslova 2 mínusy). Neplatí pro rozsahy
%	(0-9), spojovníky typu xxx-li apod.
% - Uvozovky - používat české \uv{hezké uvozovky}.
% - Zvýrazňovat \textit{kurzívou}.
% - Názvy modulů, funkcí psát \texttt{neproporcionálním fontem}.
% - URL psát do \url{}
% - In-line matematické výrazy přímo do řádky $uzavřené do dolarů$, složitější
%	na samostatnou řádku $$do dvojitých$$.
% - Pokud např. v zadání zmiňujete LATEX a TEX, použít \latex resp. \tex.

% - Pozor na předložky a spojky - nesmí se vyskytovat na konci řádku => použít
%	nedělitelnou mezeru (~).
% - Nové odstavce dvojitým odřádkováním.

% OBRÁZKY	- Vektorově (!!!) - Illustrator, free Inkscape.
% 			- Rastrově vkládat jen screeny obrazovky.
%			- V textu pak odkazovat pomocí \ref{img1}.
%			- Pro úpravu velikosti změnit scale.
% 	\begin{figure}[h]
%		\centering
%		\includegraphics[scale=1.0]{soubor.pdf/png}
%		\caption{Popisek}
%		\label{img1}
%	\end{figure}

% PŘELOŽENÍ
% - Nainstalovat doporučený MiKTeX.
% - "pdflatex <nazev-souboru>.tex"
% - Pro správné vygenerování odkazů, obsahu nutné vždy spustit 2x (!!!)
% - Pro zjednodušení stačí spouštět přibalený bat, který to udělá sám ;)

% ^^^^^^^^^^^^^^^^^^^^^^^^^^^^^^^^^^^^^^^^^^^^^^^^^^^^^^^^^^^^^^^^^^^^^^^^^^
% ^^^^^^^^^^^^^^^^^ !!!!!!!!!SMAZAT PŘI ODEVZDÁNÍ!!!!!!!!! ^^^^^^^^^^^^^^^^^
% ^^^^^^^^^^^^^^^^^^^^^^^^^^^^^^^^^^^^^^^^^^^^^^^^^^^^^^^^^^^^^^^^^^^^^^^^^^

\documentclass[
12pt,
a4paper,
pdftex,
czech,
titlepage
]{report}

\usepackage[czech]{babel}
\usepackage[utf8]{inputenc}
\usepackage{lmodern}
\usepackage{textcomp}
\usepackage[T1]{fontenc}
\usepackage{amsfonts}
\usepackage{titlesec}
\usepackage{graphicx}
\usepackage{hyperref}

\titleformat{\chapter}
  {\normalfont\LARGE\bfseries}{\thechapter}{1em}{}
\titlespacing*{\chapter}{0pt}{0ex plus 1ex minus .2ex}{2.0ex plus .2ex}
\setlength\parindent{0pt}

\begin{document}

\begin{titlepage}
	\vspace*{-2cm}
	{\centering\includegraphics[scale=1.0]{logo.pdf}\par}
	\centering
	\vspace*{2cm}
	{\Large Semestrální práce z KIV/PC\par}
	\vspace{1.5cm}
	{\Huge\bfseries Vyhledávání cest v grafu technikou DFS\par}
	\vspace{2cm}

	{\Large Petr Hlaváč\par}
	{\Large A15B0037P\par}
	{\Large petrh@students.zcu.cz\par}

	\vfill

	{\Large 1.\,12.\,2017}
\end{titlepage}

\tableofcontents
\thispagestyle{empty}
\clearpage

\chapter{Zadání}
Naprogramujte v ANSI C přenositelnou \textbf{konzolovou aplikaci}, která bude \textbf{procházet graf technikou DFS} (Depth--First--Search). Vstupem aplikace bude soubor s popisem grafu. Výstupem je pak odpovídající výčet všech cest mezi požadovanými uzly.

\medskip
Program se bude spouštět příkazem \textbf{dfs.exe} <soubor-grafu> <id1> <id2> <maxD>
\begin{itemize}
\item Symbol <soubor-grafu> zastupuje parametr - název vstupního souboru se strukturou grafu.
\item Následují identifikátory (dále jen id) dvou uzlů v grafu <id1> a <id2>, mezi kterými bude spouštěn proces hledání cest.
\item <maxD> je parametr popisující maximální délku cest, které mají být nalezeny.
\end{itemize}

Vámi vyvinutý program tedy bude vykonávat následující činnosti.
\begin{enumerate}
\item Při spuštění bez potřebných parametrů vypíše nápovědu pro jeho správné spuštění a ukončí se.
\item Při spouštění s parametry načtě zadaný vstupní soubor do vhodné struktury reprezentující graf a mezi zadanými uzly najde všechny cesty, jejichž délka nepřekročí konstatnu nastavenou vstupním parametrem.
\end{enumerate}

Váš program může být během testování spuštěn například takto:

\medskip
\framebox[1.1\width]{dfs.exe graph.csv 1 2 4} \par

\medskip
Při tomto spuštění program vyhledá všechny cesty délky 1 až 4 mezi uzly 1 a 2.

\medskip
Hotovou práci odevzdejte v jediném archivu typu ZIP prostřednictvím automatického odevzdávacího a validačního systému. Postupujte podle instrukcí uvedených na webu předmětu. Archiv nechť obsahuje všechny zdrojové soubory potřebné k přeložení programu, \textbf{makefile} pro Windows i Linux a dokumentaci ve formátu PDF vytvořenou v typografickém systému \TeX, resp. \LaTeX.

\section{Popis činnosti programu}
Program načte vstup ve formátu CSV, který vhodným zbůsobem uloží do paměti pro pozdější prohledávání. Následně spustí algoritmus, jehož výstupem bude seznam všech cest mezi dvěma uzly, které bude vypisovat na standardní výstup ve formátu popsaném dále.

\subsection{Specifikace výstupu programu}
Program bude na standardní výstup vypisovat jednotlivé cesty. Vždy na jeden řádek právě jednu cestu. Cesty budou popsány posloupností id jenotlivých uzlů oddělených pomlčkou (tj. znakem minus,  \uv{\textbf{-}}). následovat bude středník a popisky jednotlicých hran v grafu oddělené čárkou. Za posledním středníkem bude uvedena hodnota sekundární metriky - relevance cesty, podle které budou cesty seřazeny (viz Řazení výstupu).

\medskip
Například tedy pro hledání cest mezi uzly A a B:

\bigskip
\textbf{A-B};$h_{1}$;$m_{1}$

\textbf{A-F-B};$h_{2}$,$h_{3}$;$m_{2,3}$

\textbf{A-W-B};$h_{4}$,$h_{5}$;$m_{4,5}$

\textbf{A-F-O-B};$h_{2}$,$h_{6}$,$h_{7}$;$m_{2,6,7}$

\bigskip
Kde A, B, F, W, F, ... jsou popisky uzlů grafu, $h_{n}$ jsou ohodnocení hran a $m_{c1,c2...cn}$ je hodnota metrika relevance nalezené cesty.

\subsection{Řazení výstupu}
Cesty budou na standardním výstupu seřazeny podle jejich délky. Pokud bude nalezeno více cest stejné dálky, budou seřezeny podle jejich relevance následujícím způsobem. Popisky grafu nesou informaci o kalendářním datu ve formátu YYYY--MM--DD. Každá cesta budde tedy ohodnocena celým číslem, které bude odpovídat rozdílu v počtu dní mezi nejstarším a nejnovějším datem, následně budou podle toho čísla cesty se shodnou délkou seřazeny vzestupně.

\medskip
Pokud tedy nalezete cestu, kterou budou tvořit tyto 4 hrany s hodnotami:

\bigskip
2000-01-15,


2000-01-04,


2000-01-19,


2000-01-24,

\bigskip
vypočte se hodnota metriky jako počet dnů mezi daty 2000-01-04 a 2000-01-24, což je 20.

\section{Podrobné zadání}
Originální nezkrácené zadání je možné nalézt na adrese:

\medskip
\url{http://www.kiv.zcu.cz/studies/predmety/pc/doc/work/sw2017-02.pdf}

\setcounter{page}{1}

\chapter{Analýza úlohy}

\chapter{Implementace}

\chapter{Uživatelská příručka}

\chapter{Závěr}

\end{document}